
\scanpage{54}
\mspage{10}
des graphes pondérés à repères, l'opération \deleted{\unsure{s'explicite}}
se reflète par une opération correspondante au niveau de $\vec R$, par le ``gommage'' d'une partie $\vec R'$ stable par $\mathbb D_\infty$ i.e. par $\sigma_0,\sigma_1$.

2) Gommage de cercles de découpage.

On se donne maintenant une \deleted{ens. de composantes connexes de} 1-sous-variété
\inserted{(ouv. et)} \inserted{fermée} $\Sigma'$ de $\Sigma\setminus\partial X$,
correspondant à un ensemble \deleted{\unsure{A}} de \deleted{\unsure{comp.}} conn. de $\Sigma'$ i.e. une partie \inserted{$A'_0$} de $\pi_0(\unsure{\Sigma'|\Sigma})=A_0$
(arêtes ordinaires du graphe) --- et on remplace $\Sigma$ par $\Sigma\setminus\Sigma'$.
Au niveau des graphes pondérés, cela se traduit par une opération de \uline{contraction en un point} des arêtes correspondantes, d'où formation d'un nouveau graphe, et d'un \uline{morphisme de générisation} (à l'opposé de ``spécialisation'') de l'ancien vers le nouveau.
Il convient ici d'expliciter comment la pondération suit\ldots\
Au niveau des \deleted{repères} \inserted{sommets}, il est entendu que les sommets repères se trouvant sur $\Sigma'$ sont gommés en $=$ temps --- donc, sur le graphe pondéré à repères, on fait sauter l'ens des repères qui sont au dessus des arêtes $\in\unsure{A'_{0a}|A'_0}$.
\mspage{53}
